\subsection*{定理 5.6.}
\underline{设 $X$ 为满足 (LV) (5.2.2) 的光滑、不可约射影簇. 设 $\{H_\tau\}$ 为超平面线束, 其一般成员及轴分别与 $X$ 横截相交于 $X_\eta$ 和 $\Delta$. 假设投影 $\rho:\widetilde X\to\mathbb P^1$ 满足条件 (A). 则:}
\footnote{印本写作“条件 (5.3.2)”, 但 (5.3.2) 是所列出的一族伴随映射; 条件 (A) 在 5.3.5 中通过这些映射定义.}

\smallskip
\noindent
\textup{(5.6.1)}\quad
\underline{Leray 谱序列}
\[
E_2^{p,q}=H^p(\mathbb P^1,R^q\rho_*\mathbb Q_\ell)
\Longrightarrow H^{p+q}(\widetilde X,\mathbb Q_\ell)
\]
\footnote{印本的收敛项为 $H^{m+q}$; 双次数以及定理 5.6.8 中的对应陈述要求 $H^{p+q}$.}
\underline{退化.}

\noindent
\textup{(注.)}\quad
\underline{若不假设 $X$ 满足 (LV), 而仅假设一个光滑超平面截面 $Y$ 满足 (LV), 此结论仍成立.}

\subsubsection*{5.6.2.}
\underline{以 $\eta$ 表示 $\mathbb P^1$ 的泛点, $\bar\eta$ 表示几何泛点, $X_{\bar\eta}$ 表示 $\rho$ 的几何一般纤维, $\pi$ 表示 $k(\bar\eta)/k(\eta)$ 的 Galois 群.}

\underline{当 $q\neq n-1$ 时有典范同构}
\begin{align}
E_2^{0,q}
 &=H^0(\mathbb P^1,R^q\rho_*\mathbb Q_\ell)
 \simeq H^q(X_{\bar\eta},\mathbb Q_\ell),
 \tag{5.6.2.1}\\
E_2^{2,q}
 &=H^2(\mathbb P^1,R^q\rho_*\mathbb Q_\ell)
 \simeq H^q(X_{\bar\eta},\mathbb Q_\ell(-1)),
 \tag{5.6.2.2}\\
E_2^{p,q}
 &=H^p(\mathbb P^1,R^q\rho_*\mathbb Q_\ell)=0,
 \qquad p\neq0,2.
 \tag{5.6.2.3}
\end{align}
\underline{当 $q=n-1$ 时有}
\begin{align}
E_2^{0,n-1}
 &=H^0(\mathbb P^1,R^{n-1}\rho_*\mathbb Q_\ell)
 \simeq H^{n-1}(X_{\bar\eta},\mathbb Q_\ell)^\pi
 \simeq H^{n-1}(X,\mathbb Q_\ell),
 \tag{5.6.2.4}\\
E_2^{2,n-1}
 &=H^2(\mathbb P^1,R^{n-1}\rho_*\mathbb Q_\ell)
 \simeq H^{n-1}(X_{\bar\eta},\mathbb Q_\ell(-1))^\pi
 \simeq H^{n-1}(X,\mathbb Q_\ell(-1)).
 \tag{5.6.2.5}
\end{align}
\underline{(尚未给出的 $E_2^{1,n-1}$ 将在下面 5.7 中算出.)}

\underline{证明.}
由 (5.4.3), 可以应用定理 1.1 得到 (5.6.1). 前三个公式 (5.6.2.1)--(5.6.2.3) 由 $q\neq n-1$ 时 $R^q\rho_*\mathbb Q_\ell$ 为常层 (引理 5.4) 得出; 对 (5.6.2.3), 当然还要使用 $\mathbb P^1$ 的单连通性.

(5.6.2.4) 中的第一个同构由下式得出:
\[
R^{n-1}\rho_*\mathbb Q_\ell
\overset{\sim}{\longrightarrow}
\nu_*\nu^*R^{n-1}\rho_*\mathbb Q_\ell.
\]
因此
\[
H^0(\mathbb P^1,R^{n-1}\rho_*\mathbb Q_\ell)
\simeq
H^0(U,\nu^*R^{n-1}\rho_*\mathbb Q_\ell),
\]
而 $\nu^*R^{n-1}\rho_*\mathbb Q_\ell$ 是 $U$ 上的\textquotedblleft 扭常\textquotedblright{}层, 因而对应于 $\pi$-模 $H^{n-1}(X_{\bar\eta},\mathbb Q_\ell)$. (5.6.2.4) 中的第二个同构由退化性与推论 5.1.6 得出.

现在考虑 (5.6.2.5). 以 $F^\bullet$ 表示由谱序列定义的 $H^\bullet(\widetilde X,\mathbb Q_\ell)$ 上的滤过. 由退化性, 有 $F^2H^{n+1}(\widetilde X,\mathbb Q_\ell)=E_2^{2,n-1}$, 且由 (5.6.2.3), $\operatorname{gr}^1H^{n+1}(\widetilde X,\mathbb Q_\ell)=E_2^{1,n}=0$. 因而 (参见 5.1.0),
\begin{align}
E_2^{2,n-1}
 &=F^1H^{n+1}(\widetilde X,\mathbb Q_\ell) \notag\\
 &=\operatorname{Ker}\!\left(
 H^{n+1}(\widetilde X,\mathbb Q_\ell)
 \longrightarrow H^0(\mathbb P^1,R^{n+1}\rho_*\mathbb Q_\ell)
 \right) \notag\\
 &\overset{(A)}{=}
 \operatorname{Ker}\!\left(
 k^*:H^{n+1}(\widetilde X,\mathbb Q_\ell)
 \longrightarrow H^{n+1}(X_{\bar\eta},\mathbb Q_\ell)
 \right).
\tag{5.6.3}
\end{align}
由 (5.1.2), $k^*$ 可写为
\begin{equation}\tag{5.6.4}
H^{n+1}(X,\mathbb Q_\ell)
\oplus H^{n-1}(\Delta,\mathbb Q_\ell(-1))
\xrightarrow{\ m^*+h_*\ }
H^{n+1}(X_{\bar\eta},\mathbb Q_\ell).
\end{equation}
利用 (5.2.2.3), 它变为
\begin{equation}\tag{5.6.5}
\begin{split}
&\operatorname{Prim}^{n-1}(X,\mathbb Q_\ell(-1))
\oplus H^{n-3}(X,\mathbb Q_\ell(-2))
\oplus H^{n-1}(\Delta,\mathbb Q_\ell(-1))\\
&\hspace{35mm}\xrightarrow{\ m^*L+m^*L^2+h_*\ }
H^{n+1}(X_{\bar\eta},\mathbb Q_\ell).
\end{split}
\end{equation}
\footnote{印本的 (5.6.5) 略去了几个右括号, 以及第三个直和项的映射 $h_*$ 前面的加号. 所列的三个源直和项与 (5.6.4) 决定了这里给出的改正后的三项映射.}
而映射
\[
H^{n-3}(X,\mathbb Q_\ell(-2))
\xrightarrow{\ m^*L^2\ }
H^{n+1}(X_{\bar\eta},\mathbb Q_\ell)
\]
由 (LV) 为同构, 因为 $m^*L_X^2=L_{X_{\bar\eta}}^2m^*$ (参见 5.2.4 和 5.2.5). 因而 (5.6.5) 的核经由 $\operatorname{pr}_1\oplus\operatorname{pr}_3$ 对应于
\[
\operatorname{Prim}^{n-1}(X,\mathbb Q_\ell(-1))
\oplus H^{n-1}(\Delta,\mathbb Q_\ell(-1)).
\]
由 (LV) 和弱 Lefschetz 定理,
\[
H^{n-1}(\Delta,\mathbb Q_\ell(-1))
\xleftarrow[\sim]{\,L\,}
H^{n-3}(\Delta,\mathbb Q_\ell(-2))
\xleftarrow[\sim]{\ }
H^{n-3}(X,\mathbb Q_\ell(-2)).
\]
结合 (5.2.2.3), 得到
\begin{equation}\tag{5.6.6}
E_2^{2,n-1}\simeq H^{n-1}(X,\mathbb Q_\ell(-1)).
\end{equation}
(5.6.2.5) 中的另一个同构由 (5.6.2.4) 经过扭转得到.

\subsubsection*{注 5.6.7.}
Deligne 向我说明, 断言 (5.6.1) 在弱得多的假设下也能证明:
